\documentclass[11pt,a4paper]{article}
\usepackage{kotex}
\usepackage{amsmath,amssymb,booktabs,geometry}
\geometry{margin=25mm}
\title{I01 선검증 3: 비균일 위상 표지의 Aer 유한 샷 판별}
\author{RELATIONAL ATLAS}
\date{2026년 10월 4일}
\begin{document}
\maketitle

\section*{정정: 사전 삭제 검정의 확증 판정 철회}
아래 표의 구 재배분 $p$값은 raw count의 기술량으로만 남긴다.
균일 표지의 비제약 상태에서는 표지별 두 블록 확률이 달라도
가중 합은 같을 수 있는데, 구 재배분 null은 표지별 확률까지
같다는 더 강한 조건을 사용했다. 따라서 약한 합산 null에 대한
유효한 확증 검정이 아니다. 이 문제를 발견한 뒤 새 통계량과
독립 seed를 사전 고정한 선검증 4b에서 재시험한다.

\section{판정}
발전연구 주제는 여전히 정하지 않는다. Aer raw counts는
비균일/비제약 $\tau=0$에서 정확 모형이 예상한 방향의 차이를
보였다. 그러나 원래의 재배분 검정은 위의 null 오정의 때문에
확증 판정에서 제외한다. $\tau=\pi/2$의 정확 양수는 작으며,
정확한 영, 약한 양수, 유한 샷 검출은 다른 판정이다.

\section{사전 설계와 검정}
\texttt{prevalidation\_protocol3\_2026-10-04.tex}는 Aer 실행 전에
설정과 판정선을 고정했다. 이 설정은 선검증 2의 정확 확률에 근거한
검정력 파일럿으로 고른 것이므로 독립 물리 홀드아웃은 아니다.
$n=3$, $\beta=0.45$, 기준 프로젝터 폭 $\beta_r=0.15$,
$E=E_0+0.37E_0^2/3$, $E_0=\sum k_i^2$,
$\tau=0,\pi/2$, 공통 위상 $0\to0.60$ rad를 사용했다.
표지 7개를 균일하거나
$a_t=[1+0.35\cos(2\pi t/7)]/7$의 빈도로 배분하고,
비제약/$P=0$ 상태를 각각 준비했다.

축·상태·시계·위상 블록당 총 229376 shots를 사용했다.
균일 표지는 각각 32768 shots, 비균일 표지는
$(44237,39919,30216,22435,22435,30216,39918)$ shots였다.
총 336회로를 Aer statevector, CPU workers 6으로 실행했다.
두 블록의 표지별 배분은 동일했다. 원래의 분석은 표지를 합쳐
축별 성공률을 비교하되, 재배분 null은 축$\times$표지 층별
총 성공수를 조건화했다. 이 방식은 서로 다른 표지끼리 성공확률이
같다고 가정하지는 않지만, \emph{두 블록의 같은 표지 확률}이
같다고 가정한다. 합산 확률만 같은 null과 달라서, 2000회
재배분 및 사전 Bonferroni $0.00625$ 문턱은 확증 판정으로
사용하지 않는다.

\begin{center}
\begin{tabular}{@{}llrr@{}}
\toprule
표지 분포/상태 & $\tau/\pi$ & 정확 최대 확률 차이 & 구 $p$ (참고) \\
\midrule
균일/비제약 & 0 & $<10^{-15}$ & 0.63218 \\
균일/비제약 & 1/2 & $<10^{-15}$ & 0.93603 \\
균일/$P=0$ & 0 & 0 & 0.84508 \\
균일/$P=0$ & 1/2 & 0 & 0.61069 \\
비균일/비제약 & 0 & 0.002899 & 0.00050 \\
비균일/비제약 & 1/2 & 0.000497 & 0.01399 \\
비균일/$P=0$ & 0 & 0 & 0.53573 \\
비균일/$P=0$ & 1/2 & 0 & 0.79610 \\
\bottomrule
\end{tabular}
\end{center}

정확 확률 파일럿의 근사 비중심성은 이 샷 예산에서 비균일/비제약
$\tau=0$에 26.5, $\tau=\pi/2$에 2.1이었다. 따라서 뒤의
구 분석의 비검출은 정확 null의 근거가 아니다. 같은 raw counts에서 표지를
그대로 사용하는 검정은 비제약 네 조건 모두 $p=0.00050$이었고,
$P=0$ 네 조건은 검출되지 않았다. 표지 보존의 이 비교는
8개 삭제 분석 가족 외의 탐색적 진단이다.

\section{사후 구현 감사와 재가중 진단}
336개 회로의 raw count를 정확 이항확률로 표준화한 평균은
$-0.0271$, 분산은 $0.8711$이었다. 5000개 정확-이항 패널의
95\% 평균 범위 $[-0.1077,0.1074]$, 분산 범위
$[0.8208,1.2092]$ 안에 든다. 이 한 seed의 Aer 계수에서
뛰는 편향이나 분산 이상은 보이지 않았다. 모든 검출 한계가
사라졌다는 뜻은 아니다.

raw 결과를 본 \emph{뒤} 추가한 진단에서는 비균일 배분의 표지
기록을 사용해 각 표지 성공률을 동일 가중치 $1/7$로 다시 합쳤다.
그 표준화 목표에서는 정확 모형의 공통 위상 차이가 0이다.
비균일/비제약 $\tau=0$의 관측 차이에 대한 탐색적 정규근사
$p$값은 0.890, $\tau=\pi/2$에서는 0.327이었다.
이는 표지 기록과 배분 정보를 \emph{사용할 수 있다면} 표준화
분석이 가능하다는 구현 진단이지, 표지를 실제로 잃은 검출기가
그 정보를 복구했다는 뜻이 아니다. 사후 계산이므로 확증
검정으로 세지 않는다.

이번 회로는 에너지에 따른 대각 위상 유니터리를 사용했으며,
별도 시계 레지스터의 후선택을 구현하지 않았다. 따라서 이 실험만으로
일반적인 유한 시계 아래 식별 가능성을 주장할 수 없다.
또한 표지 배분은 계획한 정수로 정확히 구현했다. 무작위 표지 수의
요동, 미지 배분의 추정 오차, 기준 위상 자체의 불확실도 및 다른
cutoff는 다음 검증 게이트로 남긴다. QPU는 사용하지 않았다.

\paragraph{재현 자료.}
사전 설계: \texttt{prevalidation\_protocol3\_2026-10-04.tex}.
정확 파일럿: \texttt{../src/relational\_atlas\_i01/holdout/qiskit\_rank\_holdout3\_power.py},
\texttt{../results/qiskit\_rank\_holdout3\_power.json}.
Aer 원자료와 8검정:
\texttt{../src/relational\_atlas\_i01/holdout/qiskit\_rank\_holdout3\_shots.py},
\texttt{../results/qiskit\_rank\_holdout3\_shots.json}.
사후 감사:
\texttt{../src/relational\_atlas\_i01/holdout/qiskit\_rank\_holdout3\_audit.py},
\texttt{../results/qiskit\_rank\_holdout3\_audit.json}.
\end{document}
